\section*{Resumen del curso}
\addcontentsline{toc}{section}{Resumen del curso}

En la escala más fundamental que conocemos, hay cuatro formas en las que el contenido del universo puede interactuar: la gravitación, las interacciones nucleares fuerte y débil, y el electromagnetismo. Cada una de ellas moldea la realidad en una escala de tamaños y energías concreta. La gravitación es responsable de la formación de grandes estructuras en el universo. En el lado opuesto, las interacciones fuerte y débil tienen un alcance espacial tan reducido que solo son realmente relevantes dentro de los núcleos y en aceleradores de partículas. El electromagnetismo, por su parte, es protagonista en un rango de escalas que toca ambos regímenes. En el mundo microscópico, dicta la formación de átomos y moléculas. En el mundo macroscópico, permite la propagación de señales lumínicas hasta distancias cósmicas.

Es justo decir que el electromagnetismo es, con el permiso de la gravitación, la interacción fundamental que más directamente afecta nuestra experiencia del mundo que nos rodea. No es de extrañar, entonces, que su estudio se remonte a la Antigüedad en Egipto, quienes poseían ya conocimientos de óptica geométrica. Aun así, no ha sido hasta hace relativamente poco que tenemos descripciones completas de los fenómenos electromagnéticos; hubo que esperar hasta el s. XIX para la teoría clásica, y hasta el s. XX para la versión cuántica. 

Otra prueba tangible para los estudiantes de física de la importancia del electromagnetismo es su fuerte presencia a lo largo del grado de Física. Hay cuatro asignaturas obligatorias dedicadas al electromagnetismo: Fundamentos de Física II (Electromagnetismo 0), Electromagnetismo I y II, y Óptica, que es una especie de Electromagnetismo III. Siguiendo esta progresión, Fotónica viene a ser un Electromagnetismo IV. Sin embargo, a diferencia de las asignaturas previas, la fotónica tiene un pie firmemente puesto en el mundo de la investigación moderna. Ya no se trata meramente de aprender los fundamentos del electromagnetismo. Ahora nos centramos en adquirir técnicas y conceptos avanzados que han dado lugar a ramas de investigación actual como los metamateriales, los cristales fotónicos, la óptica no lineal o los láseres.

El objetivo principal de estas notas es adentrarnos en algunos temas importantes del campo de la fotónica. Para ello, empezamos con la \cref{sec.1-electrodinamica}, en la que discutimos los principios básicos de la electrodinámica. Ciertas partes de la explicación solapan con cursos anteriores, pero a través de la repetición de ciertas derivaciones aprovechamos para introducir nuevas herramientas matemáticas que no suelen verse antes. A continuación, la \cref{sec.2-medioslineales} describe el comportamiento de las ondas electromagnéticas cuando se propagan a través de materiales con respuesta lineal. Este formalismo nos servirá para hablar de cristales fotónicos, metamateriales y guías de onda. Después pasaremos a discutir en la \cref{sec.3-mediosnolineales} la fenomenología asociada a materiales con respuesta no lineal, mucho más interesante pero más complicada desde el punto de vista teórico. Finalmente, en la \cref{sec.4-laser} exploraremos el funcionamiento de los láseres, una herramienta primordial de la física moderna, y las ecuaciones que los describen.

Unas palabras de advertencia sobre estas notas. Hay un énfasis global en el formalismo y las derivaciones que no tiene por qué encantar a todo el mundo, al menos no en primera lectura. Aun así, la experiencia me dice que sin las matemáticas suficientes es imposible hablar, y mucho menos hacer, física ``de verdad''; al menos según mi definición que, como el profesor soy yo, es la que importa.
